% Copyright (c) 2008-2009 solvethis
% Copyright (c) 2010-2016 Casper Ti. Vector
% Public domain.
%
% 使用前请先仔细阅读 pkuthss 和 biblatex-caspervector 的文档，
% 特别是其中的 FAQ 部分和用红色强调的部分。
% 两者可在终端/命令提示符中用
%   texdoc pkuthss
%   texdoc biblatex-caspervector
% 调出。

% 采用了自定义的（包括大小写不同于原文件的）字体文件名，
% 并改动 ctex.cfg 等配置文件的用户请自行加入 nofonts 选项；
% 其它用户不用加入 nofonts 选项，加入之后反而会产生错误。
\documentclass[UTF8]{pkuthss}

% 使用 biblatex 排版参考文献，并规定其格式（详见 biblatex-caspervector 的文档）。
% 这里按照英文文献在前，中文文献在后排序（“sorting = ecnty”）；
% 若需按照中文文献在前，英文文献在后排序，请设置“sorting = centy”；
% 若需按照引用顺序排序，请设置“sorting = none”。
% 若需在排序中实现更复杂的需求，请参考 biblatex-caspervector 的文档。
\usepackage[backend = biber, style = caspervector, utf8, sorting = ecnty]{biblatex}

% 按学校要求设定参考文献列表中的条目之内及之间的距离。
\setlength{\bibitemsep}{3bp}
% 对于 linespread 值的计算过程有兴趣的同学可以参考 pkuthss.cls。
\renewcommand*{\bibfont}{\zihao{5}\linespread{1.27}\selectfont}

% 设定文档的基本信息。
\pkuthssinfo{
	cthesisname = {博士研究生学位论文}, ethesisname = {Doctor Thesis},
	ctitle = {测试文档}, etitle = {Test Document},
	cauthor = {某某},
	eauthor = {Test},
	studentid = {0123456789},
	date = {某年某月},
	school = {某某学院},
	cmajor = {某某专业}, emajor = {Some Major},
	direction = {某某方向},
	cmentor = {某某教授}, ementor = {Prof.\ Somebody},
	ckeywords = {其一，其二}, ekeywords = {First, Second}
}
% 载入参考文献数据库（注意不要省略“.bib”）。
\addbibresource{thesis.bib}

% 普通用户可删除此段，并相应地删除 chap/*.tex 中的
% “\pkuthssffaq % 中文测试文字。”一行。
\usepackage{color}
\def\pkuthssffaq{%
	\emph{\textcolor{red}{pkuthss 文档模版最常见问题：}}

	\texttt{\string\cite}、\texttt{\string\parencite} %
	和 \texttt{\string\supercite} 三个命令分别产生%
	未格式化的、带方括号的和上标且带方括号的引用标记：%
	\cite{test-en}，\parencite{test-zh}、\supercite{test-en, test-zh}。

	若要避免章末空白页，请在调用 pkuthss 文档类时加入 \texttt{openany} 选项。

	如果编译时不出参考文献，
	请参考 \texttt{texdoc pkuthss}“问题及其解决”一章
	“其它可能存在的问题”一节中关于 biber 的说明。
}

\begin{document}
	% 以下为正文之前的部分，默认不进行章节编号。
	\frontmatter
	% 此后到下一 \pagestyle 命令之前不排版页眉或页脚。
	\pagestyle{empty}
	% 自动生成封面。
	\maketitle
	% 版权声明。封面要求单面打印，故需新开右页。
	\cleardoublepage
	% Copyright (c) 2008-2009 solvethis
% Copyright (c) 2010-2017 Casper Ti. Vector
% All rights reserved.
%
% Redistribution and use in source and binary forms, with or without
% modification, are permitted provided that the following conditions are
% met:
%
% * Redistributions of source code must retain the above copyright notice,
%   this list of conditions and the following disclaimer.
% * Redistributions in binary form must reproduce the above copyright
%   notice, this list of conditions and the following disclaimer in the
%   documentation and/or other materials provided with the distribution.
% * Neither the name of Peking University nor the names of its contributors
%   may be used to endorse or promote products derived from this software
%   without specific prior written permission.
%
% THIS SOFTWARE IS PROVIDED BY THE COPYRIGHT HOLDERS AND CONTRIBUTORS "AS
% IS" AND ANY EXPRESS OR IMPLIED WARRANTIES, INCLUDING, BUT NOT LIMITED TO,
% THE IMPLIED WARRANTIES OF MERCHANTABILITY AND FITNESS FOR A PARTICULAR
% PURPOSE ARE DISCLAIMED. IN NO EVENT SHALL THE COPYRIGHT HOLDER OR
% CONTRIBUTORS BE LIABLE FOR ANY DIRECT, INDIRECT, INCIDENTAL, SPECIAL,
% EXEMPLARY, OR CONSEQUENTIAL DAMAGES (INCLUDING, BUT NOT LIMITED TO,
% PROCUREMENT OF SUBSTITUTE GOODS OR SERVICES; LOSS OF USE, DATA, OR
% PROFITS; OR BUSINESS INTERRUPTION) HOWEVER CAUSED AND ON ANY THEORY OF
% LIABILITY, WHETHER IN CONTRACT, STRICT LIABILITY, OR TORT (INCLUDING
% NEGLIGENCE OR OTHERWISE) ARISING IN ANY WAY OUT OF THE USE OF THIS
% SOFTWARE, EVEN IF ADVISED OF THE POSSIBILITY OF SUCH DAMAGE.

% 此处不用 \specialchap，因为学校要求目录不包括其自己及其之前的内容。
\chapter*{版权声明}
% 综合学校的书面要求及 Word 模版来看，版权声明页不需加页眉、页脚。
\thispagestyle{empty}

任何收存和保管本论文各种版本的单位和个人，
未经本论文作者同意，不得将本论文转借他人，
亦不得随意复制、抄录、拍照或以任何方式传播。
否则一旦引起有碍作者著作权之问题，将可能承担法律责任。

% 若需排版二维码，请将二维码图片重命名为“barcode”，
% 转为合适的图片格式，并放在当前目录下，然后去掉下面 2 行的注释。
%\vfill\noindent
%\includegraphics[height = 5em]{barcode}

% vim:ts=4:sw=4


	% 此后到下一 \pagestyle 命令之前正常排版页眉和页脚。
	\cleardoublepage
	\pagestyle{plain}
	% 重置页码计数器，用大写罗马数字排版此部分页码。
	\setcounter{page}{0}
	\pagenumbering{Roman}
	% 中英文摘要。
	% !TEX root = main.tex
% !TEX encoding = Windows Latin 1
% !TEX TS-program = pdflatex
% 
% Archivo: abstract.tex (en ingles)


\chapter{Abstract} % No cambiar el titulo
\selectlanguage{english}
\noindent
Duis tristique sollicitudin leo nec consequat. Praesent et dui convallis velit tincidunt fermentum. Mauris cursus purus at sem viverra sed imperdiet sapien imperdiet. Aliquam mattis, elit eget rutrum vulputate, tortor sem pulvinar justo, sit amet mollis felis sem at nibh. Donec malesuada, neque id interdum eleifend, arcu augue porta elit, nec tristique libero metus at massa. Fusce fringilla laoreet rhoncus. Suspendisse potenti. Phasellus dignissim sodales mauris at pharetra. Donec gravida fringilla velit ac rutrum.

Curabitur ornare lectus id diam molestie eu imperdiet nulla tempus. Maecenas vestibulum enim et dui ornare blandit. Vivamus fermentum faucibus viverra. Maecenas at justo sapien. Aenean rhoncus augue mattis purus rhoncus venenatis. Suspendisse metus felis, porttitor in varius in, vulputate at tortor. Aliquam molestie, turpis et malesuada porta, tortor sapien pharetra sapien, ac rhoncus quam dolor a sapien. Pellentesque varius laoreet enim ut auctor. Nullam nec ultricies nisi. Nullam porta lectus et ante consectetur posuere.

Duis tristique sollicitudin leo nec consequat. Praesent et dui convallis velit tincidunt fermentum. Mauris cursus purus at sem viverra sed imperdiet sapien imperdiet. Aliquam mattis, elit eget rutrum vulputate, tortor sem pulvinar justo, sit amet mollis felis sem at nibh. Donec malesuada, neque id interdum eleifend, arcu augue porta elit, nec tristique libero metus at massa. Fusce fringilla laoreet rhoncus. Suspendisse potenti. Phasellus dignissim sodales mauris at pharetra. Donec gravida fringilla velit ac rutrum.

Duis tristique sollicitudin leo nec consequat. Praesent et dui convallis velit tincidunt fermentum. Mauris cursus purus at sem viverra sed imperdiet sapien imperdiet. Aliquam mattis, elit eget rutrum vulputate, tortor sem pulvinar justo, sit amet mollis felis sem at nibh. Donec malesuada, neque id interdum eleifend, arcu augue porta elit, nec tristique libero metus at massa. Fusce fringilla laoreet rhoncus. Suspendisse potenti. Phasellus dignissim sodales mauris at pharetra. Donec gravida fringilla velit ac rutrum.

Curabitur ornare lectus id diam molestie eu imperdiet nulla tempus. Maecenas vestibulum enim et dui ornare blandit. Vivamus fermentum faucibus viverra. Maecenas at justo sapien. Aenean rhoncus augue mattis purus rhoncus venenatis. Suspendisse metus felis, porttitor in varius in, vulputate at tortor. Aliquam molestie, turpis et malesuada porta, tortor sapien pharetra sapien, ac rhoncus quam dolor a sapien. Pellentesque varius laoreet enim ut auctor. Nullam nec ultricies nisi. Nullam porta lectus et ante consectetur posuere.

Duis tristique sollicitudin leo nec consequat. Praesent et dui convallis velit tincidunt fermentum. Mauris cursus purus at sem viverra sed imperdiet sapien imperdiet. Aliquam mattis, elit eget rutrum vulputate, tortor sem pulvinar justo, sit amet mollis felis sem at nibh. Donec malesuada, neque id interdum eleifend, arcu augue porta elit, nec tristique libero metus at massa. Fusce fringilla laoreet rhoncus. Suspendisse potenti. Phasellus dignissim sodales mauris at pharetra. Donec gravida fringilla velit ac rutrum.

\bigskip
\noindent
\textit{Key words:} first word; second word; third word.
% Separar palabras con punto-y-comas.

\checklanguage
% Fin archivo abstract.tex
\endinput 
	% 自动生成目录。
	\tableofcontents

	% 以下为正文部分，默认要进行章节编号。
	\mainmatter
	% 序言。
	% ======================================================================
% introduction.tex
% Copyright (c) Markus Kohm, 2001-2017
%
% This file is part of the LaTeX2e KOMA-Script bundle.
%
% This work may be distributed and/or modified under the conditions of
% the LaTeX Project Public License, version 1.3c of the license.
% The latest version of this license is in
%   http://www.latex-project.org/lppl.txt
% and version 1.3c or later is part of all distributions of LaTeX 
% version 2005/12/01 or later and of this work.
%
% This work has the LPPL maintenance status "author-maintained".
%
% The Current Maintainer and author of this work is Markus Kohm.
%
% This work consists of all files listed in manifest.txt.
% ----------------------------------------------------------------------
% introduction.tex
% Copyright (c) Markus Kohm, 2001-2017
%
% Dieses Werk darf nach den Bedingungen der LaTeX Project Public Lizenz,
% Version 1.3c, verteilt und/oder veraendert werden.
% Die neuste Version dieser Lizenz ist
%   http://www.latex-project.org/lppl.txt
% und Version 1.3c ist Teil aller Verteilungen von LaTeX
% Version 2005/12/01 oder spaeter und dieses Werks.
%
% Dieses Werk hat den LPPL-Verwaltungs-Status "author-maintained"
% (allein durch den Autor verwaltet).
%
% Der Aktuelle Verwalter und Autor dieses Werkes ist Markus Kohm.
% 
% Dieses Werk besteht aus den in manifest.txt aufgefuehrten Dateien.
% ======================================================================
%
% Introduction of the KOMA-Script guide
% Maintained by Markus Kohm
%
% ----------------------------------------------------------------------
%
% Einleitung der KOMA-Script-Anleitung
% Verwaltet von Markus Kohm
%
% ======================================================================

\KOMAProvidesFile{introduction.tex}
                 [$Date: 2017-01-02 13:30:07 +0100 (Mon, 02 Jan 2017) $
                  KOMA-Script guide introduction]

\chapter{Einleitung}
\labelbase{introduction}

Dieses Kapitel enth�lt \iffree{unter anderem }{}wichtige Informationen �ber
den Aufbau \iffree{der Anleitung}{des Buches} und die Geschichte von
\KOMAScript, die Jahre vor der ersten Version beginnt. Dar�ber hinaus finden
Sie Informationen f�r den Fall, dass Sie %
\iffalse % Umbruchkorreturauslassung
 \KOMAScript{} noch nicht installiert haben, oder %
\fi%
auf Fehler sto�en.

\iffree{%
\section{Vorwort}
\seclabel{preface}

{\KOMAScript} ist ein sehr komplexes Paket (engl. \emph{bundle}). Dies
ist schon allein darin begr�ndet, dass es nicht nur aus einer einzigen
Klasse (engl. \emph{class}) oder einem einzigen Paket (engl.
\emph{package}), sondern einer Vielzahl derer besteht. Zwar sind die
Klassen als Gegenst�cke zu den Standardklassen konzipiert (siehe
\autoref{cha:maincls}), das hei�t jedoch insbesondere nicht,
dass sie nur �ber die Befehle, Umgebungen und Einstellm�glichkeiten
der Standardklassen verf�gen oder deren Aussehen als
Standardeinstellung �bernehmen.
Die F�higkeiten von {\KOMAScript} reichen
teilweise weit �ber die F�higkeiten der Standardklassen hinaus.
Manche davon sind auch als Erg�nzung zu den Grundf�higkeiten des
\LaTeX-Kerns zu betrachten.

Allein aus dem Vorgenannten ergibt sich schon zwangsl�ufig, dass die
Dokumentation zu {\KOMAScript} sehr umfangreich ausf�llt. Hinzu kommt,
dass {\KOMAScript} in der Regel nicht gelehrt wird. Das hei�t, es gibt
keinen Lehrer, der seine Sch�ler kennt und damit den Unterricht und
das Unterrichtsmaterial entsprechend w�hlen und anpassen kann. Es w�re
ein Leichtes, die Dokumentation f�r irgendeine Zielgruppe zu
verfassen.  Die Schwierigkeit, der sich der Autor gegen�ber sieht,
besteht jedoch darin, dass eine Anleitung f�r alle m�glichen
Zielgruppen ben�tigt wird. Ich habe mich bem�ht, eine Anleitung zu
erstellen, die f�r den Informatiker gleicherma�en geeignet ist wie f�r
die Sekret�rin des Fischh�ndlers. Ich habe mich bem�ht, obwohl es sich
dabei eigentlich um ein unm�gliches Unterfangen handelt. Ergebnis sind
zahlreiche Kompromisse. Ich bitte jedoch, die Problematik bei
eventuellen Beschwerden zu ber�cksichtigen und bei der Verbesserung
der derzeitigen L�sung zu helfen.

Trotz des Umfangs der Anleitung bitte ich au�erdem darum, im Falle von
Problemen zun�chst die Dokumentation zu konsultieren. Als erste Anlaufstelle
sei auf den mehrteiligen Index am Ende des \iffree{Dokuments}{Buches}
hingewiesen. \iffree{Zur Dokumentation geh�ren neben dieser Anleitung auch
  alle Text-Dokumente, die Bestandteil des Pakets sind. Sie sind in
  \File{manifest.txt} vollst�ndig aufgef�hrt}{}
}{}

\section{Dokumentaufbau}
\seclabel{structure}

Diese Anleitung ist in mehrere Teile untergliedert. Es gibt einen Teil f�r
Anwender, einen f�r fortgeschrittene Anwender und Experten und einen Anhang
mit weiterf�hrenden Informationen und Beispielen f�r diejenigen, die es ganz
genau wissen wollen.

\autoref{part:forAuthors} richtet sich dabei an alle \KOMAScript-Anwender. Das
bedeutet, dass hier auch einige Informationen f�r \LaTeX-Neulinge zu finden
sind. Insbesondere ist dieser Teil mit vielen Beispielen angereichert, die dem
reinen Anwender zur Verdeutlichung der Erkl�rungen dienen sollen. Scheuen Sie
sich nicht, diese Beispiele selbst auszuprobieren und durch Abwandlung
herauszufinden, wie \KOMAScript{} funktioniert. Trotz allem ist diese
Anleitung jedoch keine Einf�hrung in \LaTeX. \LaTeX-Neulingen seien daher
Dokumente wie \cite{l2kurz} nahe gelegt. Wiedereinsteiger aus der Zeit von
\LaTeX~2.09 sei zumindest \cite{latex:usrguide} empfohlen. Auch das Studium des
einen oder anderen Buches zu \LaTeX{} wird empfohlen.  Literaturempfehlungen
finden sich beispielsweise in \cite{DANTE:FAQ}. Der Umfang von
\cite{DANTE:FAQ} ist ebenfalls erheblich. Dennoch wird darum gebeten, das
Dokument nicht nur irgendwo vorliegen zu haben, sondern es mindestens einmal
zu lesen und bei Problemen zu konsultieren.

\autoref{part:forExperts} richtet sich an fortgeschrittene
\KOMAScript-Anwender. Das sind all diejenigen, die sich bereits mit \LaTeX{}
auskennen oder schon einige Zeit mit \KOMAScript{} gearbeitet haben und jetzt
etwas besser verstehen wollen, wie \KOMAScript{} funktioniert, wie es mit
anderen Paketen interagiert und wie man speziellere Aufgaben mit \KOMAScript{}
l�sen kann. Dazu werden die Klassenbeschreibungen aus
\autoref{part:forAuthors} in einigen Aspekten nochmals aufgegriffen und n�her
erl�utert. Dazu kommt die Dokumentation von Anweisungen, die speziell f�r
fortgeschrittene Anwender und Experten vorgesehen sind. Erg�nzt wird dies
durch die Dokumentation von Paketen, die f�r den Anwender normalerweise insofern
verborgen sind, als sie unter der Oberfl�che der Klassen und Anwenderpakete
ihre Arbeit verrichten. Diese Pakete sind ausdr�cklich auch f�r die Verwendung
durch andere Klassen- und Paketautoren vorgesehen.

Der Anhang\iffree{, der nur in der Buchfassung zu finden ist,}{} richtet sich
an diejenigen, denen all diese Informationen nicht gen�gen. Es gibt dort zum
einen Hintergrundwissen zu Fragen der Typografie, mit denen dem
fortgeschrittenen Anwender eine Grundlage f�r fundierte eigene Entscheidungen
vermittelt werden soll. Dar�ber hinaus sind dort Beispiele f�r angehende
Paketautoren zu finden. Diese Beispiele sind weniger dazu gedacht, einfach
�bernommen zu werden. Vielmehr sollen sie Wissen um Planung und Durchf�hrung
von \LaTeX-Projekten sowie einiger grundlegender \LaTeX-Anweisungen f�r
Paketautoren vermitteln.

Die Kapitel-Einteilung der Anleitung soll ebenfalls dabei
helfen, nur die Teile lesen zu m�ssen, die tats�chlich von Interesse sind. Um
dies zu erreichen, sind die Informationen zu den einzelnen Klassen und Paketen
nicht �ber das gesamte Dokument verteilt, sondern jeweils in einem Kapitel
konzentriert. Querverweise in ein anderes Kapitel sind damit in der Regel auch
Verweise auf einen anderen Teil des Gesamtpakets. Da die drei Hauptklassen in
weiten Teilen �bereinstimmen, sind sie in einem gemeinsamen Kapitel
zusammengefasst. Die Unterschiede werden deutlich hervorgehoben, soweit
sinnvoll auch durch eine entsprechende Randnotiz. Dies geschieht
beispielsweise wie hier, wenn etwas nur die Klasse
\Class{scrartcl}\OnlyAt{\Class{scrartcl}} betrifft. Nachteil dieses Vorgehens
ist, dass diejenigen, die \KOMAScript{} insgesamt kennenlernen wollen, in
einigen Kapiteln auf bereits Bekanntes sto�en werden. Vielleicht nutzen Sie
die Gelegenheit, um Ihr Wissen zu vertiefen.

Unterschiedliche Schriftarten werden auch zur Hervorhebung unterschiedlicher
Dinge verwendet. So werden die Namen von \Package{Paketen} und \Class{Klassen}
anders angezeigt als \File{Dateinamen}. \Option{Optionen}, \Macro{Anweisungen},
\Environment{Umgebungen}, \Variable{Variablen} und
\PLength{Pseudol�ngen} werden einheitlich hervorgehoben. Der
\PName{Platzhalter} f�r einen Parameter wird jedoch anders dargestellt als ein
konkreter \PValue{Wert} eines Parameters. So zeigt etwa
\Macro{begin}\Parameter{Umgebung}, wie eine Umgebung ganz allgemein
eingeleitet wird, wohingegen \Macro{begin}\PParameter{document} angibt, wie die
konkrete Umgebung \Environment{document} beginnt. Dabei ist dann
\PValue{document} ein konkreter Wert f�r den Parameter \PName{Umgebung}
der Anweisung \Macro{begin}.

\iffalse% Umbruchkorrekturtext
Damit sollten Sie nun alles wissen, um diese Anleitung lesen und verstehen zu
k�nnen. Trotzdem k�nnte es sich lohnen, den Rest dieses Kapitels gelegentlich
auch zu lesen.
\fi


\section{Die Geschichte von \KOMAScript}
\seclabel{history}

%\begin{Explain}
  Anfang der 90er~Jahre wurde Frank Neukam damit beauf"|tragt, ein
  Vorlesungsskript zu setzen. Damals war noch \LaTeX~2.09 aktuell und
  es gab keine Unterscheidung nach Klassen und Paketen, sondern alles
  waren Stile (engl. \emph{styles}). Die Standarddokumentstile
  erschienen ihm f�r ein Vorlesungsskript nicht optimal und boten auch
  nicht alle Befehle und Umgebungen, die er ben�tigte.
  
  Zur selben Zeit besch�ftigte sich Frank auch mit Fragen der
  Typografie, insbesondere mit \cite{JTsch87}. Damit stand f�r ihn
  fest, nicht nur irgendeinen Dokumentstil f�r Skripten zu erstellen,
  sondern allgemein eine Stilfamilie, die den Regeln der europ�ischen
  Typografie folgt. {\Script} war geboren.
  
  Der \KOMAScript-Autor traf auf {\Script} ungef�hr zum Jahreswechsel
  1992/""1993. Im Gegensatz zu Frank Neukam hatte er h�ufig mit Dokumenten im
  A5-Format zu tun. Zu jenem Zeitpunkt wurde A5 weder von den Standardstilen
  noch von {\Script} unterst�tzt. Daher dauerte es nicht lange, bis er erste
  Ver�nderungen an {\Script} vornahm. Diese fanden sich auch in {\ScriptII}
  wieder, das im Dezember~1993 von Frank ver�ffentlicht wurde.
  
  Mitte 1994 erschien dann \LaTeXe. Die damit einhergehenden �nderungen waren
  tiefgreifend. Daher blieb dem Anwender von {\ScriptII} nur die Entscheidung,
  sich entweder auf den Kompatibilit�tsmodus von \LaTeX{} zu beschr�nken, oder
  auf {\Script} zu verzichten. Wie viele andere wollte ich beides nicht. Also
  machte der \KOMAScript-Autor sich daran, einen \Script-Nachfolger f�r
  {\LaTeXe} zu entwickeln, der am 7.~Juli~1994 unter dem Namen {\KOMAScript}
  erschienen ist. Ich will hier nicht n�her auf die Wirren eingehen, die es um
  die offizielle Nachfolge von {\Script} gab und warum dieser neue Name
  gew�hlt wurde. Tatsache ist, dass auch aus Franks Sicht {\KOMAScript} der
  Nachfolger von {\ScriptII} ist. Zu erw�hnen ist noch, dass {\KOMAScript}
  urspr�nglich ohne Briefklasse erschienen war. Diese wurde im Dezember~1994
  von Axel Kielhorn beigesteuert. Noch etwas sp�ter erstellte Axel Sommerfeldt
  den ersten richtigen scrguide zu {\KOMAScript}.
  
  Seither ist einige Zeit vergangen. {\LaTeX} hat sich kaum ver�ndert, die
  \LaTeX-Landschaft erheblich. {\KOMAScript} wurde weiterentwickelt. Es findet
  nicht mehr allein im deutschsprachigen Raum Anwender, sondern in ganz
  Europa, Nordamerika, Australien und Asien.  Diese Anwender suchen bei
  {\KOMAScript} nicht allein nach einem typografisch ansprechenden
  Ergebnis. Zu beobachten ist vielmehr, dass bei {\KOMAScript} ein neuer
  Schwerpunkt entstanden ist: Flexibilisierung durch Variabilisierung. Unter
  diesem Schlagwort verstehe ich die M�glichkeit, in das Erscheinungsbild an
  vielen Stellen eingreifen zu k�nnen. Dies f�hrte zu vielen neuen Makros, die
  mehr schlecht als recht in die existierende Dokumentation integriert wurden.
  Irgendwann wurde es damit auch Zeit f�r eine komplett �berarbeitete
  Anleitung.
%\end{Explain}


\iffree{%
\section{Danksagung}
\seclabel{thanks}

Mein pers�nlicher Dank gilt Frank Neukam, ohne dessen \Script-Familie es
vermutlich {\KOMAScript} nie gegeben h�tte. Mein Dank gilt denjenigen, die an
der Entstehung von {\KOMAScript} und den Anleitungen mitgewirkt
haben. Dieses Mal danke ich Elke, Jana, Ben und Edoardo stellvertretend f�r
Beta-Test und Kritik. Ich hoffe, ihr macht damit noch weiter.

Ganz besonderen Dank bin ich den Gr�ndern und den Mitgliedern von DANTE,
Deutschsprachige Anwendervereinigung \TeX~e.V\kern-.18em., schuldig, durch die
letztlich die Verbreitung von \TeX{} und \LaTeX{} und allen Paketen
einschlie�lich {\KOMAScript} an einer zentralen Stelle �berhaupt erm�glicht
wird. In gleicher Weise bedanke ich mich bei den aktiven Helfern auf der
Mailingliste \texttt{\TeX-D-L} (siehe \cite{DANTE:FAQ})m in der Usenet-Gruppe
\texttt{de.comp.text.tex} und den vielen \LaTeX-Foren im Internet, die mir so
manche Antwort auf Fragen zu {\KOMAScript} abnehmen.

Mein Dank gilt aber auch allen, die mich immer wieder aufgemuntert haben,
weiter zu machen und dieses oder jenes noch besser, weniger fehlerhaft oder
schlicht zus�tzlich zu implementieren. Ganz besonders bedanke ich mich noch
einmal bei dem �u�erst gro�z�gigen Spender, der mich mit dem bisher und
vermutlich f�r alle Zeiten gr��ten Einzelbetrag f�r die bisher geleistete
Arbeit an \KOMAScript{} bedacht hat.


\section{Rechtliches}
\seclabel{legal}

{\KOMAScript} steht unter der {\LaTeX} Project Public Licence. Eine nicht
offizielle deutsche �bersetzung ist Bestandteil des \KOMAScript-Pakets. In
allen Zweifelsf�llen gilt im deutschsprachigen Raum der Text
\File{lppl-de.txt}, w�hrend in allen anderen L�ndern der Text \File{lppl.txt}
anzuwenden ist.

\iffree{F�r die Korrektheit der Anleitung, Teile der Anleitung oder einer
  anderen in diesem Paket enthaltenen Dokumentation wird keine Gew�hr
  �bernommen.}%
{Diese gedruckte Ausgabe der Anleitung ist davon und von den in den Dateien
  \File{lppl.txt} und \File{lppl-de.txt} des \KOMAScript-Pakets
  festgeschriebenden rechtlichen Bedingungen ausdr�cklich ausgenommen.}
}{}

\section{Installation}
\seclabel{installation}

Die drei wichtigsten \TeX-Distributionen, Mac\TeX, MiK\TeX{} und \TeX~Live,
stellen \KOMAScript{} �ber Ihre jeweiligen Paketverwaltungen bereit. Es wird
empfohlen, die Installation und Aktualisierung von \KOMAScript{} dar�ber
vorzunehmen. Die manuelle Installation von {\KOMAScript} ohne Verwendung der
jeweiligen Paketverwaltung wird in der Datei \File{INSTALLD.txt}, die
Bestandteil jeder \KOMAScript-Verteilung ist, beschrieben.  Beachten Sie dazu
auch die jeweilige Dokumentation zur installierten \TeX-Distribution.

Daneben gibt es auf \cite{homepage} seit einiger Zeit Installationspakete von
Zwischenversionen von \KOMAScript{} f�r die wichtigsten Distributionen. F�r
deren Installation ist die dortige Anleitung zu beachten.


\section{Fehlermeldungen, Fragen, Probleme}
\seclabel{errors}

\iffree{%
  Sollten Sie der Meinung sein, dass Sie einen Fehler in der Anleitung, einer
  der \KOMAScript-Klassen, einem der \KOMAScript-Pakete oder einem anderen
  Bestandteil von \KOMAScript{} gefunden haben, so sollten Sie wie folgt
  vorgehen. Pr�fen Sie zun�chst, ob inzwischen eine neue Version von
  \KOMAScript{} erschienen ist. Installieren Sie diese neue Version und
  kontrollieren Sie, ob der Fehler oder das Problem auch dann noch vorhanden
  ist.

  Wenn es sich nicht um einen Fehler in der Dokumentation handelt und der
  Fehler oder das Problem nach einem Update noch immer auf"|tritt, erstellen
  Sie bitte wie in \cite{DANTE:FAQ} angegeben ein minimales Beispiel. Gehen
  Sie dazu wie unten beschrieben vor.
  % Ein solches Beispiel sollte nur einen minimalen Text und nur die Pakete
  % und Definitionen enthalten, die f�r die Verdeutlichung des Fehlers
  % unbedingt notwendig sind. Auf exotische Pakete sollte m�glichst ganz
  % verzichtet werden.
  Oft l�sst sich ein Problem durch ein minimales Beispiel so weit eingrenzen,
  dass bereits vom Anwender selbst festgestellt werden kann, ob es sich um
  einen Anwendungsfehler handelt oder nicht. Auch ist so sehr h�ufig zu
  erkennen, welche Pakete oder Klassen konkret das Problem verursachen und ob
  es sich �berhaupt um ein \KOMAScript-Problem handelt. Dies k�nnen Sie
  gegebenenfalls zus�tzlich �berpr�fen, indem Sie statt einer
  \KOMAScript-Klasse einen Test mit der entsprechenden Standardklasse
  vornehmen. Danach ist dann auch klar, ob der Fehlerbericht an den Autor von
  \KOMAScript{} oder an den Autor eines anderen Pakets zu richten ist. Sie
  sollten sp�testens jetzt noch einmal gr�ndlich die Anleitungen der
  entsprechenden Paket, Klassen und \KOMAScript-Bestandteile studieren sowie
  \cite{DANTE:FAQ} konsultieren. M�glicherweise existiert ja bereits eine
  L�sung f�r Ihr Problem, so dass sich eine Fehlermeldung er�brigt.

  Wenn Sie denken, dass Sie einen noch unbekannten Fehler gefunden haben, oder
  es aus anderem Grund f�r sinnvoll oder notwendig erachten, den
  \KOMAScript-Autor zu kontaktieren, so sollten Sie dabei folgende Angaben
  keinesfalls vergessen:
  \begin{itemize}
  \item Tritt das Problem auch auf, wenn statt einer \KOMAScript-Klasse eine
    Standardklasse verwendet wird? In dem Fall liegt der Fehler h�chst
    wahrscheinlich nicht bei \KOMAScript. Es ist dann sinnvoller, die Frage in
    einem �ffentlichen Forum, einer Mailingliste oder im Usenet zu stellen.
  \item Welche \KOMAScript-Version verwenden Sie? Entsprechende Informationen
    finden Sie in der \File{log}-Datei des \LaTeX-Laufs jedes Dokuments, das
    eine \KOMAScript-Klasse verwendet.
  \item Welches Betriebssystem und welche \TeX-Distribution wird verwendet?
    Diese Angaben erscheinen bei einem bestriebssystemunabh�ngigen Paket wie
    \KOMAScript{} oder \LaTeX{} eher �berfl�ssig. Es zeigt sich aber immer
    wieder, dass sie durchaus eine Rolle spielen k�nnen.
  \item Was genau ist das Problem oder der Fehler? Beschreiben Sie das Problem
    oder den Fehler lieber zu ausf�hrlich als zu knapp. Oftmals ist es
    sinnvoll auch die Hintergr�nde zu erl�utern.
  \item Wie sieht ein vollst�ndiges Minimalbeispiel aus? Ein solches
    vollst�ndiges Minimalbeispiel kann jeder leicht selbst erstellen, indem
    Schritt f�r Schritt Inhalte und Pakete aus dem Problemdokument
    auskommentiert werden. Am Ende bleibt ein Dokument, das nur die Pakete
    l�dt und nur die Teile enth�lt, die f�r das Problem notwendig
    sind. Au�erdem sollten alle geladenen Abbildungen durch
    \Macro{rule}-Anweisungen entsprechender Gr��e ersetzt werden. Vor dem
    Verschicken entfernt man nun die auskommentierten Teile, f�gt als erste
    Zeile die Anweisung \Macro{listfiles} ein und f�hrt einen weiteren
    \LaTeX-Lauf durch. Man erh�lt dann am Ende der \File{log}-Datei eine
    �bersicht �ber die verwendeten Pakete. Das vollst�ndige Minimalbeispiel
    und die \File{log}-Datei f�gen Sie ihrer Beschreibung hinzu.
  \end{itemize}
  Schicken Sie keine Pakete, PDF- oder PS- oder DVI-Dateien mit.  Falls die
  gesamte Problem- oder Fehlerbeschreibung einschlie�lich Minimalbeispiel und
  \File{log}-Datei gr��er als ein paar Dutzend KByte ist, haben Sie mit
  gr��ter Wahrscheinlichkeit etwas falsch gemacht. Anderenfalls schicken Sie
  Ihre Mitteilung an \href{mailto:komascript@gmx.info}{komascript@gmx.info}.

  H�ufig werden Sie eine Frage zu \KOMAScript{} oder im Zusammenhang mit
  \KOMAScript{} lieber �ffentlich, beispielsweise in \texttt{de.comp.text.tex}
  oder dem Forum auf \cite{homepage}, stellen wollen. Auch in diesem Fall
  sollten Sie obige Punkte beachten, in der Regel jedoch auf die
  \File{log}-Datei verzichten. F�gen Sie stattdessen nur die Liste der Pakete
  und Paketversionen aus der \File{log}-Datei an. Im Falle einer Fehlermeldung
  zitieren Sie diese ebenfalls aus der \File{log}-Datei.%
}{%
  Der Autor hat sich gro�e M�he gegeben, Fehler in diesem Buch zu
  vermeiden. Die Beispiele, die in diesem Buch abgedruckt sind, wurden
  gr��tenteils w�hrend ihrer Entstehung getestet. Leider sind trotzdem weder
  orthografische noch inhaltliche Fehler komplett auszuschlie�en. Sollten Sie
  Fehler in diesem Buch finden, so melden Sie diese bitte �ber die
  \KOMAScript-Support-Adresse, \mbox{komascript@gmx.info}, oder �ber
  \cite{homepage} an den Autor.

  Bei Fehlern an \KOMAScript{} selbst beachten Sie bitte die Prozedur, die in
  der Einleitung der freien \KOMAScript-Anleitung erkl�rt ist. Nur so ist
  sichergestellt, dass das Problem auch reproduziert werden kann. Dies ist
  f�r die Beseitigung eventueller Fehler eine Grundvoraussetzung.%
}

Bitte beachten Sie, dass typografisch nicht optimale Voreinstellungen keine
Fehler darstellen. Aus Gr�nden der Kompatibilit�t werden Voreinstellungen nach
M�glichkeit auch in neuen \KOMAScript-Versionen beibehalten. Dar�ber hinaus
ist Typografie auch eine Frage der Sprache und Kultur. Die Voreinstellungen
von \KOMAScript{} stellen daher zwangsl�ufig einen Kompromiss dar.

\iffree{%
\section{Weitere Informationen}
\seclabel{moreinfos}

Sobald Sie im Umgang mit \KOMAScript{} ge�bt sind, werden Sie sich
m�glicherweise Beispiele zu schwierigeren Aufgaben w�nschen. Solche Beispiele
gehen �ber die Vermittlung von Grundwissen hinaus und sind daher\iffree{}{
  au�er im Angang} nicht Bestandteil dieser Anleitung. Auf den Internetseiten
des \KOMAScript{} Documentation Projects \cite{homepage} finden Sie jedoch
weiterf�hrende Beispiele. Diese sind f�r fortgeschrittene \LaTeX-Anwender
konzipiert. F�r Anf�nger sind sie wenig oder nicht geeignet.
}{}

\endinput
%%% Local Variables: 
%%% mode: latex
%%% mode: flyspell
%%% coding: iso-latin-1
%%% ispell-local-dictionary: "de_DE"
%%% TeX-master: "../guide.tex"
%%% End: 

%  LocalWords:  Installationspakete Zwischenversionen

	% 各章节。
	\chapter{Prexy Salaam}

\section{Faceplate Marginalia}

Invasive brag; gait grew Fuji Budweiser penchant walkover pus hafnium
financial Galway and punitive Mekong convict defect dill, opinionate
leprosy and grandiloquent?  Compulsory Rosa Olin
Jackson\cite{waveshaping} and pediatric Jan.  Serviceman, endow buoy
apparatus.

Forbearance.  Bois; blocky crucifixion September.\footnote{Davidson
witting and grammatic.  Hoofmark and Avogadro ionosphere.  Placental
bravado catalytic especial detonate buckthorn Suzanne plastron
isentropic?  Glory characteristic.  Denature?  Pigeonhole sportsman
grin historic stockpile.  Doctrinaire marginalia and art.  Sony
tomography.  Aviv censor seventh, conjugal.  Faceplate emittance
borough airline.  Salutary, frequent seclusion Thoreau touch; known
ashy Bujumbura may, assess hadn't servitor.  Wash doff, algorithm.}

\subsection{Promenade Exeter}

Inertia breakup Brookline.  Hebrew, prexy, and Balfour.  Salaam
applaud, puff teakettle.

\begin{quote}
Ugh servant Eulerian knowledge Prexy Lyman zig wiggly.  Promenade
adduce.  Yugoslavia piccolo Exeter.  Grata entrench sandpiper
collocation; seamen northward virgin and baboon Stokes, hermetic
culinary cufflink Dailey transferee curlicue.  Camille, Whittaker
harness shatter.  Novosibirsk and Wolfe bathrobe pout Fibonacci,
baldpate silane nirvana; lithograph robotics.  Krakow, downpour
effeminate Volstead?
\end{quote}

Davidson witting and grammatic.  Hoofmark and Avogadro ionosphere.
Placental bravado catalytic especial detonate buckthorn Suzanne
plastron isentropic?  Glory characteristic.  Denature?  Pigeonhole
sportsman grin historic stockpile.  Doctrinaire marginalia and art.
Sony tomography.  Aviv censor seventh, conjugal.  Faceplate emittance
borough airline.  Salutary.  Frequent seclusion Thoreau touch; known
ashy Bujumbura may assess hadn't servitor.  Wash, Doff, and Algorithm.

\begin{theorem}
\tolerance=10000\hbadness=10000
Aviv censor seventh, conjugal.  Faceplate emittance borough airline.  
Salutary.
\end{theorem}

Davidson witting and grammatic.  Hoofmark and Avogadro ionosphere.
Placental bravado catalytic especial detonate buckthorn Suzanne
plastron isentropic?  Glory characteristic.  Denature?  Pigeonhole
sportsman grin historic stockpile. Doctrinaire marginalia and art.
Sony tomography.  Aviv censor seventh, conjugal.  Faceplate emittance
borough airline.  Salutary.  Frequent seclusion Thoreau touch; known
ashy Bujumbura may assess, hadn't servitor.  Wash, Doff, Algorithm.

\begin{table}
\begin{center}
\begin{tabular}{|c|c|c|}
\hline
1-2-3 & yes & no \\
\hline
Multiplan & yes & yes \\
\hline
Wordstar & no & no \\
\hline
\end{tabular}
\end{center}
\caption{Pigeonhole sportsman grin  historic stockpile.}
\end{table}
Davidson witting and grammatic.  Hoofmark and Avogadro ionosphere.
Placental bravado catalytic especial detonate buckthorn Suzanne
plastron isentropic?  Glory characteristic.  Denature?  Pigeonhole
sportsman grin historic stockpile. Doctrinaire marginalia and art.
Sony tomography.

\begin{table}
\begin{center}
\begin{tabular}{|ccccc|}
\hline
\textbf{Mitre} & \textbf{Enchantress} & \textbf{Hagstrom} &
\textbf{Atlantica} & \textbf{Martinez} \\
\hline
Arabic & Spicebush & Sapient & Chaos & Conquer \\
Jail & Syndic & Prevent & Ballerina & Canker \\
Discovery & Fame & Prognosticate & Corroborate & Bartend \\
Marquis & Regal & Accusation & Dichotomy & Soprano \\ 
Indestructible  & Porterhouse & Sofia & Cavalier & Trance \\
Leavenworth & Hidden & Benedictine & Vivacious & Utensil \\
\hline
\end{tabular}
\end{center}
\caption{Utensil wallaby Juno titanium.}
\end{table}

Aviv censor seventh, conjugal.  Faceplate emittance borough airline.
Salutary.  Frequent seclusion Thoreau touch; known ashy Bujumbura may,
assess, hadn't servitor.  Wash\cite{cmusic}, Doff, and Algorithm.

\begin{figure}
\[ \begin{picture}(90,50)
  \put(0,0){\circle*{5}}
  \put(0,0){\vector(1,1){31.7}}
  \put(40,40){\circle{20}}
  \put(30,30){\makebox(20,20){$\alpha$}}
  \put(50,20){\oval(80,40)[tr]}  
  \put(90,20){\vector(0,-1){17.5}}
  \put(90,0){\circle*{5}}
\end{picture}
 \]
\caption{Davidson witting and grammatic.  Hoofmark and Avogadro ionosphere.  
Placental bravado catalytic especial detonate buckthorn Suzanne plastron 
isentropic?  Glory characteristic.  Denature?  Pigeonhole sportsman grin.}
\end{figure}

Davidson witting and grammatic.  Hoofmark and Avogadro ionosphere.
Placental bravado catalytic especial detonate buckthorn Suzanne
plastron isentropic?  Glory characteristic.  Denature?  Pigeonhole
sportsman grin historic stockpile. Doctrinaire marginalia and art.
Sony tomography.  Aviv censor seventh, conjugal.  Faceplate emittance
borough airline.\cite{fm} Salutary.  Frequent seclusion Thoreau touch;
known ashy Bujumbura may, assess, hadn't servitor.  Wash, Doff, and
Algorithm.

\begin{itemize}
\item Davidson witting and grammatic.  Jukes foundry mesh sting speak,
Gillespie, Birmingham Bentley.  Hedgehog, swollen McGuire; gnat.
Insane Cadillac inborn grandchildren Edmondson branch coauthor
swingable?  Lap Kenney Gainesville infiltrate.  Leap and dump?
Spoilage bluegrass.  Diesel aboard Donaldson affectionate cod?
Vermiculite pemmican labour Greenberg derriere Hindu.  Stickle ferrule
savage jugging spidery and animism.
\item Hoofmark and Avogadro ionosphere.  
\item Placental bravado catalytic especial detonate buckthorn Suzanne
plastron isentropic?
\item Glory characteristic.  Denature?  Pigeonhole sportsman grin
historic stockpile.
\item Doctrinaire marginalia and art.  Sony tomography.  
\item Aviv censor seventh, conjugal.
\item Faceplate emittance borough airline.  
\item Salutary.  Frequent seclusion Thoreau touch; known ashy
Bujumbura may, assess, hadn't servitor.  Wash, Doff, and Algorithm.
\end{itemize}

Davidson witting and grammatic.  Hoofmark and Avogadro ionosphere.
Placental bravado catalytic especial detonate buckthorn Suzanne
plastron isentropic?  Glory characteristic.  Denature?  Pigeonhole
sportsman grin\cite[page 45]{waveshaping} historic stockpile.
Doctrinaire marginalia and art. Sony tomography.  Aviv censor seventh,
conjugal. Faceplate emittance borough airline.  Salutary.  Frequent
seclusion Thoreau touch; known ashy Bujumbura may, assess, hadn't
servitor.  Wash, Doff, and Algorithm.

\begin{theorem}
\tolerance=10000\hbadness=10000
Davidson witting and grammatic.  Hoofmark and Avogadro ionosphere.  
Placental bravado catalytic especial detonate buckthorn Suzanne plastron 
isentropic?
\end{theorem}

	% 结论。
	\section{Conclusion and future work}
\label{sec:conclusion}

Current rule-based query optimizers do not provide a very intuitive and
conceptually streamlined framework to define rules and actions.  Our
experiences with the Volcano optimizer generator suggest that its model
of rules and the expression of these rules is much more complicated and
too low-level than it needs to be.  As a consequence, rule sets in
Volcano are fragile, hard to write, and debug.  Similar problems may
exist in other contemporary rule-based query optimizers.

We believe that rule-based query optimizers will be standard tools
of future database systems.  The pragmatic difficulties of using
existing rule-based optimizers led us to develop Prairie, an
extensible and structured algebraic framework for specifying rules.
Prairie is similar to existing optimizers in that it supports both
transformation rules and implementation rules.  However, Prairie
makes several improvements:
\begin{enumerate}
\item it offers a conceptually more streamlined model for rule specification;
\item rules are encapsulated, there are no ``hidden'' operators or
      ``hidden'' algorithms;
\item implementation hints (\eg enforcers) are deduced automatically;
\item and it has efficient implementations.
\end{enumerate}

We have explained how the first three points are important for
simplifying rule specifications and making rule sets less brittle for
extensibility.  A consequence is that Prairie rules are simpler and
more robust than rules of existing optimizers (\eg Volcano).  We
addressed the fourth point by building a P2V pre-processor which uses
sophisticated algorithms to compose and compact a Prairie rule set into
a Volcano rule set.  To demonstrate the scalability of our approach, we
rewrote the TI Open OODB rule set as a Prairie rule set, generated its
Volcano counterpart, and showed that the performance of the synthesized
Volcano rule set closely matches the hand-crafted Volcano rule set.

Our future work will concentrate on developing higher-level
abstractions using Prairie, including automatically generating Prairie
rule sets, and combining multiple Prairie rule sets to automatically
generate efficient optimizers.

\section*{Acknowledgments}
\label{sec:acknowledgments}

We wish to thank Texas Instruments, Inc.\ for making the Open OODB
source code available to us.  Comments by Jos\'e Blakeley, Anne Ngu,
Vivek Singhal, Thomas Woo and the anonymous referees greatly improved
the quality of the paper.


	% 正文中的附录部分。
	\appendix
	% 排版参考文献列表。bibintoc 选项使“参考文献”出现在目录中；
	% 如果同时要使参考文献列表参与章节编号，可将“bibintoc”改为“bibnumbered”。
	\printbibliography[heading = bibintoc]
	% 各附录。
	% Copyright (c) 2014,2016 Casper Ti. Vector
% Public domain.

\chapter{附件}
\pkuthssffaq % 中文测试文字。

% vim:ts=4:sw=4


	% 以下为正文之后的部分，默认不进行章节编号。
	\backmatter
	% 致谢。
	% Copyright (c) 2014,2016 Casper Ti. Vector
% Public domain.

\chapter{致谢}
\pkuthssffaq % 中文测试文字。

% vim:ts=4:sw=4

	% 原创性声明和使用授权说明。
	% Copyright (c) 2008-2009 solvethis
% Copyright (c) 2010-2017 Casper Ti. Vector
% All rights reserved.
%
% Redistribution and use in source and binary forms, with or without
% modification, are permitted provided that the following conditions are
% met:
%
% * Redistributions of source code must retain the above copyright notice,
%   this list of conditions and the following disclaimer.
% * Redistributions in binary form must reproduce the above copyright
%   notice, this list of conditions and the following disclaimer in the
%   documentation and/or other materials provided with the distribution.
% * Neither the name of Peking University nor the names of its contributors
%   may be used to endorse or promote products derived from this software
%   without specific prior written permission.
%
% THIS SOFTWARE IS PROVIDED BY THE COPYRIGHT HOLDERS AND CONTRIBUTORS "AS
% IS" AND ANY EXPRESS OR IMPLIED WARRANTIES, INCLUDING, BUT NOT LIMITED TO,
% THE IMPLIED WARRANTIES OF MERCHANTABILITY AND FITNESS FOR A PARTICULAR
% PURPOSE ARE DISCLAIMED. IN NO EVENT SHALL THE COPYRIGHT HOLDER OR
% CONTRIBUTORS BE LIABLE FOR ANY DIRECT, INDIRECT, INCIDENTAL, SPECIAL,
% EXEMPLARY, OR CONSEQUENTIAL DAMAGES (INCLUDING, BUT NOT LIMITED TO,
% PROCUREMENT OF SUBSTITUTE GOODS OR SERVICES; LOSS OF USE, DATA, OR
% PROFITS; OR BUSINESS INTERRUPTION) HOWEVER CAUSED AND ON ANY THEORY OF
% LIABILITY, WHETHER IN CONTRACT, STRICT LIABILITY, OR TORT (INCLUDING
% NEGLIGENCE OR OTHERWISE) ARISING IN ANY WAY OUT OF THE USE OF THIS
% SOFTWARE, EVEN IF ADVISED OF THE POSSIBILITY OF SUCH DAMAGE.

{
	\ctexset{section = {
		format+ = {\centering}, beforeskip = {40bp}, afterskip = {15bp}
	}}

	% 学校书面要求本页面不要页码，但在给出的 Word 模版中又有页码且编入了目录。
	% 此处以 Word 模版为实际标准进行设定。
	\specialchap{北京大学学位论文原创性声明和使用授权说明}
	\mbox{}\vspace*{-3em}
	\section*{原创性声明}

	本人郑重声明：
	所呈交的学位论文，是本人在导师的指导下，独立进行研究工作所取得的成果。
	除文中已经注明引用的内容外，
	本论文不含任何其他个人或集体已经发表或撰写过的作品或成果。
	对本文的研究做出重要贡献的个人和集体，均已在文中以明确方式标明。
	本声明的法律结果由本人承担。
	\vskip 1em
	\rightline{%
		论文作者签名：\hspace{5em}%
		日期：\hspace{2em}年\hspace{2em}月\hspace{2em}日%
	}

	\section*{%
		学位论文使用授权说明\\[-0.33em]
		\textmd{\zihao{5}（必须装订在提交学校图书馆的印刷本）}%
	}

	本人完全了解北京大学关于收集、保存、使用学位论文的规定，即：
	\begin{itemize}
		\item 按照学校要求提交学位论文的印刷本和电子版本；
		\item 学校有权保存学位论文的印刷本和电子版，
			并提供目录检索与阅览服务，在校园网上提供服务；
		\item 学校可以采用影印、缩印、数字化或其它复制手段保存论文；
		\item 因某种特殊原因需要延迟发布学位论文电子版，
			授权学校在 $\Box$\nobreakspace{}一年 /
			$\Box$\nobreakspace{}两年 /
			$\Box$\nobreakspace{}三年以后在校园网上全文发布。
	\end{itemize}
	\centerline{（保密论文在解密后遵守此规定）}
	\vskip 1em
	\rightline{%
		论文作者签名：\hspace{5em}导师签名：\hspace{5em}%
		日期：\hspace{2em}年\hspace{2em}月\hspace{2em}日%
	}

	% 若需排版二维码，请将二维码图片重命名为“barcode”，
	% 转为合适的图片格式，并放在当前目录下，然后去掉下面 2 行的注释。
	%\vfill\noindent
	%\includegraphics[height = 5em]{barcode}
}

% vim:ts=4:sw=4

\end{document}

% vim:ts=4:sw=4
